\begin{pastexam}[label={px:ntu-108-ee-dsb-12}]{臺大電機所 108 學年度（B）第 12 題}
（複選；每題 10 分，每個選項單獨計分，答對 2 分，答錯倒扣 2 分）

Which of the following statements are true?
\begin{enumerate}[label=(\Alph*).]
\item Performing preorder traversal on BST produces a sorted list.
\item The $k$-th smallest element in a BST with $n$ elements can be found with $O(n)$ in time.
\item Every subtree of a BST is another BST.
\item The minimal AVL tree with height 5 has 20 nodes. (The height of a tree with one node is 0.)
\item When the structure of an AVL tree changes (e.g., with insertion or deletion), one single
rotation is sufficient to restore the AVL tree property.
\end{enumerate}
\end{pastexam}

\begin{pastsoln}
答案是 (B)(C)(D)。 (A)(B)(C) 考二元搜尋樹（第~9~章），(D)(E) 考 AVL 樹（第~10~章）。

\textbf{(A) 錯}。 排好的順序來自\textbf{中序}走訪（資結書定理~9.1）。 前序先訪根：以 $2$ 為根、
左右子節點 $1$、$3$ 的 BST，前序得 $2, 1, 3$，不是排好的。

\textbf{(B) 對}。 做中序走訪，數到第 $k$ 個訪問的節點就是第 $k$ 小（資結書定理~9.1）。 中序
走訪每個節點最多經過常數次，整棵走完是 $O(n)$，提早停下只會更少，所以 $O(n)$ 一定做得到。 樹
不平衡時高度可以是 $n - 1$，因此不能說 $O(\lg n)$；若每個節點另存子樹大小，可以做到 $O(h)$
（資結書第~9.4~節、虛擬碼~9.8），$h \le n - 1$，也在 $O(n)$ 之內。

\textbf{(C) 對}。 BST 的性質是：每個節點 $x$，左子樹的鍵都小於 $x$，右子樹的鍵都大於 $x$。
子樹裡的每個節點在原樹中就滿足這個條件，而且它的左右子樹在子樹裡和在原樹裡完全一樣，所以子樹
本身也是 BST。

\textbf{(D) 對}。 設 $N(h)$ 是高度 $h$ 的 AVL 樹最少的節點數。 要節點最少，根的一邊高度
$h - 1$、另一邊 $h - 2$，兩邊也都取最少，所以
\[
N(h) = N(h-1) + N(h-2) + 1,\qquad N(0) = 1,\ N(1) = 2
\]
（資結書第~10.2.2~節）。 依序算：$N(2) = 2 + 1 + 1 = 4$，$N(3) = 4 + 2 + 1 = 7$，
$N(4) = 7 + 4 + 1 = 12$，$N(5) = 12 + 7 + 1 = 20$。

\textbf{(E) 錯}。 兩個理由。 第一，插入時的 LR、RL 情況要做\textbf{雙旋轉}（兩次單旋轉），一次
單旋轉修不好（資結書第~10.2.3~節）。 例如依序插入 $3, 1, 2$：根 $3$ 失衡，是 LR 情況；只對 $3$
右旋，得到根 $1$、右子節點 $3$、$3$ 的左子節點 $2$，根 $1$ 仍然失衡。 要先對 $1$ 左旋、再對 $3$
右旋，得到根 $2$。 第二，刪除之後，旋轉可能讓子樹高度減少 $1$，失衡會一路往上傳，最壞要在
路徑上每一層各旋轉一次，共 $O(\lg n)$ 次（資結書第~10.2.6~節）。
\end{pastsoln}
